\documentclass[a4paper,10pt]{article}
\usepackage{lamia-pedagogy}

\begin{document}

% =========================================================================
% PAGE 1 : Compétence 01 - Compréhension et communication orales
% =========================================================================
\ficheHeader{1 Heure}{\makebox[2.5cm]{\dotfill}}{1\textsuperscript{ère} AM}{2026 / 2027}

\vspace{-3pt}
\competenceBanner{\faHeadphones\ Compétence 01 : Compréhension et communication orales}

\begin{minipage}[t]{0.48\textwidth}
\begin{tcolorbox}[
  enhanced,
  colback=secondarylight,
  colframe=secondary,
  arc=1.5mm,leftrule=3pt,rightrule=0.5pt,toprule=0.5pt,bottomrule=0.5pt,
  top=3pt,bottom=3pt,left=6pt,right=6pt,
  title={\small\textbf{\faBullseye\ Objectifs d'apprentissage}},
  coltitle=secondary!90!black,attach boxed title to top left={yshift=-2mm,xshift=4mm},
  boxed title style={colback=white,boxrule=0.5pt,colframe=secondary,arc=1mm,top=1pt,bottom=1pt}
]
\vspace{1pt}
{\small
\begin{itemize}[leftmargin=*,noitemsep,topsep=1pt]
  \item[\faCheckCircle] Identifier le thème de la situation de communication.
  \item[\faCheckCircle] Identifier les unités de sens.
  \item[\faCheckCircle] S'exprimer en fonction de la situation de communication.
\end{itemize}
}
\end{tcolorbox}
\end{minipage}
\hfill
\begin{minipage}[t]{0.48\textwidth}
\begin{tcolorbox}[
  enhanced,
  colback=accentlight,
  colframe=accent,
  arc=1.5mm,leftrule=3pt,rightrule=0.5pt,toprule=0.5pt,bottomrule=0.5pt,
  top=3pt,bottom=3pt,left=6pt,right=6pt,
  title={\small\textbf{\faExclamationTriangle\ Difficultés constatées}},
  coltitle=accent!90!black,attach boxed title to top left={yshift=-2mm,xshift=4mm},
  boxed title style={colback=white,boxrule=0.5pt,colframe=accent,arc=1mm,top=1pt,bottom=1pt}
]
\vspace{1pt}
{\small
\begin{itemize}[leftmargin=*,noitemsep,topsep=1pt]
  \item[\faTimesCircle] Difficulté à repérer le thème lors de l'écoute du support.
  \item[\faTimesCircle] Confusion entre thème et informations secondaires.
\end{itemize}
}
\end{tcolorbox}
\end{minipage}

\vspace{1pt}
\begin{tcolorbox}[
  enhanced,
  colback=white,
  colframe=cardborder,
  arc=1.5mm,
  top=2pt,bottom=2pt,left=6pt,right=6pt,
  title={\scriptsize\textbf{\faUsers\ Élèves concernés par la remédiation}},
  coltitle=primary,attach boxed title to top left={yshift=-1.8mm,xshift=3mm},
  boxed title style={colback=primarylight,boxrule=0.4pt,colframe=cardborder,arc=1mm,top=0.5pt,bottom=0.5pt}
]
\vspace{1pt}
{\scriptsize
\begin{tabularx}{\linewidth}{@{}l X l X@{}}
1. & \makebox[\linewidth]{\dotfill} & 6. & \makebox[\linewidth]{\dotfill} \\
2. & \makebox[\linewidth]{\dotfill} & 7. & \makebox[\linewidth]{\dotfill} \\
3. & \makebox[\linewidth]{\dotfill} & 8. & \makebox[\linewidth]{\dotfill} \\
4. & \makebox[\linewidth]{\dotfill} & 9. & \makebox[\linewidth]{\dotfill} \\
5. & \makebox[\linewidth]{\dotfill} & 10. & \makebox[\linewidth]{\dotfill} \\
\end{tabularx}
}
\end{tcolorbox}

\vspace{1pt}
\begin{tcolorbox}[
  enhanced,
  colback=secondary!5,
  colframe=secondary!40,
  arc=1.5mm,
  top=3pt,bottom=3pt,left=7pt,right=7pt
]
\textbf{\faCompass\ Déroulement de l'activité (Mise en situation) :}
\begin{itemize}[leftmargin=*,noitemsep,topsep=1pt]
  \item[\faCommentDots] Avez-vous déjà visité un zoo avec vos parents ?
  \item[\faCommentDots] Si oui, dites ce qu'on trouve dedans ?
\end{itemize}
\end{tcolorbox}

\vspace{1pt}
\begin{tcolorbox}[
  enhanced,
  colback=amber!5!white,
  colframe=primary!70,
  arc=2mm,
  top=3pt,bottom=3pt,left=8pt,right=8pt,
  title={\small\textbf{\faVolumeUp\ Support audio transcrit : \guillemotleft\ La gazelle \guillemotright}},
  coltitle=primary,attach boxed title to top left={yshift=-2mm,xshift=4mm},
  boxed title style={colback=white,boxrule=0.5pt,colframe=primary!70,arc=1mm,top=1pt,bottom=1pt}
]
\vspace{1pt}
{\small\itshape
La gazelle est un mammifère sauvage qui vit dans le désert. Elle est herbivore, elle se nourrit d'herbes, de plantes et de feuilles d'arbres. Son poil est de couleur rousse sur le dos et blanche sur le ventre. Elle a une queue noire et quatre pattes fines terminées par des sabots. Sa petite tête porte deux cornes. La gazelle pèse jusqu'à 30 kilos.
\begin{center}
\textbf{La gazelle est un animal très rapide.}
\end{center}
}
\end{tcolorbox}

\vspace{3pt}
{\small
\setlength{\tabcolsep}{2.5pt}

% 1ère écoute
\noindent\badge[primary]{\faStepForward\ 1\textsuperscript{ère} écoute (Pré-écoute)}

\vspace{1pt}
\noindent
\begin{tabular}{@{}p{5.7cm} p{4.3cm} p{4.1cm} p{3.8cm}@{}}
$\bullet$ \textbf{Qui parle dans cette vidéo ?} & \qcmCircle a) un adulte & \qcmCircle b) un enfant & \\
$\bullet$ \textbf{De quoi parle-t-il ?} & \qcmCircle a) du lion & \qcmCircle b) de la gazelle & \qcmCircle c) du fennec \\
$\bullet$ \textbf{La gazelle est un animal :} & \qcmCircle a) sauvage & \qcmCircle b) domestique & \\
\end{tabular}

\vspace{3pt}
% 2ème écoute
\noindent\badge[primary]{\faPlay\ 2\textsuperscript{ème} écoute (Écoute active)}

\vspace{1pt}
\noindent
\begin{tabular}{@{}p{5.7cm} p{4.3cm} p{4.1cm} p{3.8cm}@{}}
$\bullet$ \textbf{La gazelle est herbivore, elle se nourrit :} & \qcmCircle a) d'herbe & \qcmCircle b) de viande & \\
$\bullet$ \textbf{Où vit-elle ?} & \qcmCircle a) dans la savane africaine & \qcmCircle b) dans le jardin & \\
$\bullet$ \textbf{Combien pèse la gazelle ?} & \multicolumn{3}{l@{}}{\textbf{Réponse :} \makebox[11.5cm]{\dotfill}} \\
\end{tabular}

\vspace{3pt}
% 3ème écoute
\noindent\badge[primary]{\faCheckDouble\ 3\textsuperscript{ème} écoute (Post-écoute)}

\vspace{1pt}
\noindent
\begin{tabular}{@{}p{5.7cm} p{4.3cm} p{4.1cm} p{3.8cm}@{}}
$\bullet$ \textbf{Les pattes de la gazelle sont-elles :} & \qcmCircle a) fines & \qcmCircle b) grosses & \\
$\bullet$ \textbf{Sa queue est :} & \qcmCircle a) noire & \qcmCircle b) jaune & \\
\multicolumn{4}{@{}l@{}}{$\bullet$ \textbf{Comment est décrite la gazelle à la fin de la vidéo ?} \quad \textbf{Réponse :} \makebox[8.5cm]{\dotfill}} \\
\end{tabular}
}

\vspace{4pt}
\begin{tcolorbox}[
  enhanced,
  colback=secondarylight,
  colframe=secondary,
  arc=2mm,
  leftrule=3.5pt,rightrule=0.6pt,toprule=0.6pt,bottomrule=0.6pt,
  top=6pt,bottom=6pt,left=8pt,right=8pt
]
\begin{minipage}[c]{0.58\textwidth}
\textbf{\faPenSquare\ Synthèse : Complète l'énoncé suivant}\\[4pt]
{\small
\noindent La gazelle est un animal \dotfill ,\\[2pt]
elle est \dotfill ,\\[2pt]
elle pèse environ \dotfill ,\\[2pt]
elle vit dans \dotfill .
}
\end{minipage}
\hfill
\begin{minipage}[c]{0.38\textwidth}
\centering
\includegraphics[width=\linewidth,height=3.8cm,keepaspectratio]{images/gazelle_p1.png}
\end{minipage}
\end{tcolorbox}

\newpage
% =========================================================================
% PAGE 2 : Fiche élève à découper (2 colonnes A5)
% =========================================================================

\begin{tikzpicture}[remember picture, overlay]
\draw[dashed,thick,gray!60] ([yshift=-1cm]current page.north) -- ([yshift=1cm]current page.south);
\node[fill=white,rotate=90,inner sep=3pt] at (current page.center) {\textcolor{gray!70}{\faCut\ \scriptsize\textit{Ligne de découpe centrale}}};
\end{tikzpicture}

\newcommand{\modernGazelleCard}{%
\begin{minipage}[t][27.6cm][t]{0.96\linewidth}
\begin{tcolorbox}[
  enhanced,
  colback=primarylight,
  colframe=primary,
  arc=2mm,
  top=3pt,bottom=3pt,left=6pt,right=6pt
]
\textbf{\faFile*\ Fiche Élève : La gazelle} \hfill {\scriptsize\badge[primary]{1\textsuperscript{ère} AM}}\\[2pt]
{\scriptsize
\faUsers\ \textbf{Nom du groupe :} \makebox[4.5cm]{\dotfill}\\[2.5pt]
\textbf{Membres du groupe (4 élèves) :}\\[1.5pt]
\begin{tabularx}{\linewidth}{@{}l X l X@{}}
\textbf{1.} & \makebox[\linewidth]{\dotfill} & \textbf{2.} & \makebox[\linewidth]{\dotfill} \\
\textbf{3.} & \makebox[\linewidth]{\dotfill} & \textbf{4.} & \makebox[\linewidth]{\dotfill} \\
\end{tabularx}
}
\end{tcolorbox}

\vspace{1pt}
\begin{center}
\badge[primary]{\faQuestionCircle\ Questionnaire de compréhension orale}
\end{center}

\vspace{1pt}
{\scriptsize
\textbf{\faStepForward\ 1\textsuperscript{ère} écoute :}\\
\textbf{1/} Qui parle dans cette vidéo ?
\cardOptTwo{un adulte}{un enfant}

\textbf{2/} De quoi parle-t-il ?
\cardOptThree{du lion}{de la gazelle}{du fennec}

\textbf{3/} La gazelle est un animal :
\cardOptTwo{sauvage}{domestique}

\vspace{2pt}
\textbf{\faPlay\ 2\textsuperscript{ème} écoute :}\\
\textbf{1/} La gazelle est herbivore, elle se nourrit :
\cardOptTwo{d'herbe}{de la viande}

\textbf{2/} Où vit la gazelle ?
\cardOptTwo{dans la savane}{dans le jardin}

\textbf{3/} Combien pèse la gazelle ?\par\vspace{1pt}
\noindent\makebox[\linewidth][l]{\hspace*{2mm}\textbf{Rép. :} \dotfill}\par\vspace{2pt}

\vspace{2pt}
\textbf{\faCheckDouble\ 3\textsuperscript{ème} écoute :}\\
\textbf{1/} Les pattes de la gazelle sont-elles :
\cardOptTwo{fines}{grosses}

\textbf{2/} Sa queue est :
\cardOptTwo{noire}{jaune}

\textbf{3/} Comment est décrite la gazelle à la fin de la vidéo ?\par\vspace{1pt}
\noindent\makebox[\linewidth][l]{\hspace*{2mm}\textbf{Rép. :} \dotfill}\par\vspace{2pt}
}

\vspace{2pt}
\begin{tcolorbox}[
  enhanced,
  colback=secondarylight,
  colframe=secondary,
  arc=1.5mm,
  top=2.5pt,bottom=2.5pt,left=5pt,right=5pt,
  title={\scriptsize\textbf{\faPen\ Synthèse : Complète l'énoncé suivant}},
  coltitle=secondary!90!black,attach boxed title to top left={yshift=-1.8mm,xshift=3mm},
  boxed title style={colback=white,boxrule=0.4pt,colframe=secondary,arc=1mm,top=0.5pt,bottom=0.5pt}
]
\vspace{1pt}
{\scriptsize
\noindent La gazelle est un animal \dotfill ,\\[2pt]
elle est \dotfill ,\\[2pt]
elle se nourrit de \dotfill ,\\[2pt]
elle vit dans la \dotfill .
}
\end{tcolorbox}

\vspace{1pt}
\begin{center}
\includegraphics[width=0.98\linewidth,height=7.2cm,keepaspectratio]{images/gazelle_schema.png}
\end{center}
\end{minipage}%
}

\noindent
\begin{minipage}[t]{0.485\textwidth}
\modernGazelleCard
\end{minipage}
\hfill
\begin{minipage}[t]{0.485\textwidth}
\modernGazelleCard
\end{minipage}

\newpage
% =========================================================================
% PAGE 3 : Compétence 02 - Lecture
% =========================================================================
\ficheHeader{1 Heure}{\makebox[2.5cm]{\dotfill}}{1\textsuperscript{ère} AM}{2026 / 2027}

\vspace{-2pt}
\competenceBanner{\faBookOpen\ Compétence 02 : Lecture et phonétique}

\begin{minipage}[t]{0.48\textwidth}
\begin{tcolorbox}[
  enhanced,
  colback=secondarylight,
  colframe=secondary,
  arc=1.5mm,leftrule=3pt,rightrule=0.5pt,toprule=0.5pt,bottomrule=0.5pt,
  top=3pt,bottom=3pt,left=6pt,right=6pt,
  title={\small\textbf{\faBullseye\ Critères \& Objectifs}},
  coltitle=secondary!90!black,attach boxed title to top left={yshift=-2mm,xshift=4mm},
  boxed title style={colback=white,boxrule=0.5pt,colframe=secondary,arc=1mm,top=1pt,bottom=1pt}
]
\vspace{1pt}
{\scriptsize
\begin{itemize}[leftmargin=*,noitemsep,topsep=1pt]
  \item[\faCheck] Activer la correspondance graphie/phonie.
  \item[\faCheck] Réaliser une lecture fluide et expressive.
  \item[\faCheck] Respecter l'intonation des différents types de phrases.
  \item[\faCheck] Distinguer les sons : \textbf{/g/ -- /3/}, \textbf{/s/ -- /z/}, \textbf{/k/ -- /s/}.
\end{itemize}
}
\end{tcolorbox}
\end{minipage}
\hfill
\begin{minipage}[t]{0.48\textwidth}
\begin{tcolorbox}[
  enhanced,
  colback=accentlight,
  colframe=accent,
  arc=1.5mm,leftrule=3pt,rightrule=0.5pt,toprule=0.5pt,bottomrule=0.5pt,
  top=3pt,bottom=3pt,left=6pt,right=6pt,
  title={\small\textbf{\faExclamationTriangle\ Obstacles \& Difficultés}},
  coltitle=accent!90!black,attach boxed title to top left={yshift=-2mm,xshift=4mm},
  boxed title style={colback=white,boxrule=0.5pt,colframe=accent,arc=1mm,top=1pt,bottom=1pt}
]
\vspace{1pt}
{\scriptsize
\begin{itemize}[leftmargin=*,noitemsep,topsep=1pt]
  \item[\faTimesCircle] Difficulté à associer la lettre au phonème correspondant.
  \item[\faTimesCircle] Décodage laborieux gênant la fluence.
  \item[\faTimesCircle] Difficulté à moduler l'intonation (expressivité).
\end{itemize}
}
\end{tcolorbox}
\end{minipage}

\vspace{1pt}
\begin{tcolorbox}[
  enhanced,
  colback=secondarylight,
  colframe=secondary,
  arc=1.5mm,leftrule=3.5pt,rightrule=0.6pt,toprule=0.6pt,bottomrule=0.6pt,
  top=2.5pt,bottom=2.5pt,left=8pt,right=8pt
]
{\small
\textbf{\faUsers\ Modalité :} Travail en groupes \hfill \textbf{\faChalkboard\ Support :} Tableau, cahier
}
\end{tcolorbox}

\vspace{1pt}
\begin{tcolorbox}[
  enhanced,
  colback=white,
  colframe=cardborder,
  arc=1.5mm,
  top=2pt,bottom=2pt,left=6pt,right=6pt,
  title={\scriptsize\textbf{\faUsers\ Élèves ayant des difficultés de lecture}},
  coltitle=primary,attach boxed title to top left={yshift=-1.8mm,xshift=3mm},
  boxed title style={colback=primarylight,boxrule=0.4pt,colframe=cardborder,arc=1mm,top=0.5pt,bottom=0.5pt}
]
\vspace{1pt}
{\scriptsize
\begin{tabularx}{\linewidth}{@{}l X l X@{}}
1. & \makebox[\linewidth]{\dotfill} & 6. & \makebox[\linewidth]{\dotfill} \\
2. & \makebox[\linewidth]{\dotfill} & 7. & \makebox[\linewidth]{\dotfill} \\
3. & \makebox[\linewidth]{\dotfill} & 8. & \makebox[\linewidth]{\dotfill} \\
4. & \makebox[\linewidth]{\dotfill} & 9. & \makebox[\linewidth]{\dotfill} \\
5. & \makebox[\linewidth]{\dotfill} & 10. & \makebox[\linewidth]{\dotfill} \\
\end{tabularx}
}
\end{tcolorbox}

\vspace{1pt}
\begin{tcolorbox}[
  enhanced,
  colback=primarylight,
  colframe=primary,
  arc=1.5mm,
  top=2.5pt,bottom=2.5pt,left=6pt,right=6pt,
  title={\small\textbf{\faChalkboardTeacher\ Mise en {\oe}uvre de la remédiation}},
  coltitle=white,colbacktitle=primary
]
\vspace{1pt}
{\small
\begin{tabularx}{\linewidth}{|X|X|}
\hline
\rowcolor{secondarylight}
\textbf{\faUserTie\ Rôle de l'enseignante} & \textbf{\faUserGraduate\ Rôle des apprenants} \\
\hline
$\bullet$ Présente les activités et le matériel. & $\bullet$ Lisent silencieusement puis à haute voix. \\
$\bullet$ Explicite clairement les consignes. & $\bullet$ Écoutent et posent des questions. \\
$\bullet$ Oriente, régule et encourage. & $\bullet$ S'entraident et échangent en binômes. \\
\hline
\end{tabularx}
}
\end{tcolorbox}

\vspace{1pt}
\begin{tcolorbox}[
  enhanced,
  colback=accentlight,
  colframe=accent,
  arc=1.5mm,
  top=2.5pt,bottom=2.5pt,left=7pt,right=7pt
]
\textbf{\faLightbulb\ Moment de mise en confiance -- Devinette : Qui suis-je ?}\\[1pt]
{\small\itshape
\guillemotleft\ Je travaille à l'hôpital, je soigne les malades, j'utilise un stéthoscope. Qui suis-je ?\ \guillemotright\ \hfill \textbf{Réponse :} \makebox[3cm]{\dotfill}
}
\end{tcolorbox}

\vspace{1pt}
{\small
\badge[primary]{\faTasks\ Activités d'entraînement phonétique et de lecture}

\vspace{2pt}
\textbf{Activité 01 (Fluence de lecture) :} \textit{Lis attentivement le passage suivant :}\\
\begin{tcolorbox}[enhanced,colback=white,colframe=cardborder,arc=1mm,top=3pt,bottom=3pt,left=6pt,right=6pt]
\textit{\guillemotleft\ Chaque jour, Amira va à l'école à 8 heures du matin. Dans la classe, elle retrouve ses camarades et écoute sa maîtresse. Elle apprend à lire et à écrire, elle aime dessiner les fleurs et les animaux.\ \guillemotright}
\end{tcolorbox}

\vspace{2pt}
\textbf{Activité 02 (Discrimination de la lettre \guillemotleft\ C \guillemotright) :} \textit{Classe ces mots dans le tableau :}\\
\centerline{\textbf{Carotte -- cerise -- coq -- bicyclette -- cube -- citron}}
\vspace{2pt}
\begin{tabularx}{\linewidth}{|X|X|}
\hline
\rowcolor{secondarylight}
\centering\textbf{La lettre \guillemotleft\ C \guillemotright\ se prononce /k/} & \centering\textbf{La lettre \guillemotleft\ C \guillemotright\ se prononce /s/} \tabularnewline
\hline
\rule{0pt}{1.8em} & \tabularnewline
\hline
\end{tabularx}

\vspace{3pt}
\textbf{Activité 03 (Les enveloppes sonores) :} \textit{Prépare deux enveloppes : une pour le son \textbf{/g/} et l'autre pour \textbf{/3/}.}\\
\textit{Mets chaque mot dans l'enveloppe adéquate :} \textbf{Singe -- gaspiller -- guide -- garçon -- plage -- gomme.}

\vspace{3pt}
\textbf{Activité 04 (Intonation expressive) :} \textit{Lis ces phrases en respectant la ponctuation et l'intonation :}\\
$\bullet$ Les enfants visitent le zoo. \hfill $\bullet$ Aimez-vous les animaux ? \hfill $\bullet$ Quel beau métier ! \hspace{2cm}\null
}

\newpage
% =========================================================================
% PAGE 4 : Compétence 03 - Compréhension de l'écrit (Phase 1)
% =========================================================================
\ficheHeader{1 Heure}{\makebox[2.5cm]{\dotfill}}{1\textsuperscript{ère} AM}{2026 / 2027}

\vspace{-2pt}
\competenceBanner{\faBookReader\ Compétence 03 : Compréhension de l'écrit}

\begin{minipage}[t]{0.48\textwidth}
\begin{tcolorbox}[
  enhanced,
  colback=secondarylight,
  colframe=secondary,
  arc=1.5mm,leftrule=3pt,rightrule=0.5pt,toprule=0.5pt,bottomrule=0.5pt,
  top=3pt,bottom=3pt,left=6pt,right=6pt,
  title={\small\textbf{\faBullseye\ Critères 1 et 2}},
  coltitle=secondary!90!black,attach boxed title to top left={yshift=-2mm,xshift=4mm},
  boxed title style={colback=white,boxrule=0.5pt,colframe=secondary,arc=1mm,top=1pt,bottom=1pt}
]
\vspace{1pt}
{\scriptsize
Identifier le thème général et le champ lexical relatif au texte.
\begin{itemize}[leftmargin=*,noitemsep,topsep=1pt]
  \item[\faCheckCircle] Repérer les indices textuels thématiques.
  \item[\faCheckCircle] Sélectionner le vocabulaire lié au domaine.
\end{itemize}
}
\end{tcolorbox}
\end{minipage}
\hfill
\begin{minipage}[t]{0.48\textwidth}
\begin{tcolorbox}[
  enhanced,
  colback=accentlight,
  colframe=accent,
  arc=1.5mm,leftrule=3pt,rightrule=0.5pt,toprule=0.5pt,bottomrule=0.5pt,
  top=3pt,bottom=3pt,left=6pt,right=6pt,
  title={\small\textbf{\faExclamationTriangle\ Difficultés identifiées}},
  coltitle=accent!90!black,attach boxed title to top left={yshift=-2mm,xshift=4mm},
  boxed title style={colback=white,boxrule=0.5pt,colframe=accent,arc=1mm,top=1pt,bottom=1pt}
]
\vspace{1pt}
{\scriptsize
\begin{itemize}[leftmargin=*,noitemsep,topsep=1pt]
  \item[\faTimesCircle] L'élève peine à repérer le thème central.
  \item[\faTimesCircle] Difficulté à regrouper les mots d'un même champ lexical.
\end{itemize}
}
\end{tcolorbox}
\end{minipage}

\vspace{1pt}
\begin{tcolorbox}[
  enhanced,
  colback=secondarylight,
  colframe=secondary,
  arc=1.5mm,leftrule=3.5pt,rightrule=0.6pt,toprule=0.6pt,bottomrule=0.6pt,
  top=2.5pt,bottom=2.5pt,left=8pt,right=8pt
]
{\small
\textbf{\faUsers\ Modalité de travail :} En groupe puis individuel.
}
\end{tcolorbox}

\vspace{1pt}
\begin{tcolorbox}[
  enhanced,
  colback=white,
  colframe=cardborder,
  arc=1.5mm,
  top=2pt,bottom=2pt,left=6pt,right=6pt,
  title={\scriptsize\textbf{\faUsers\ Liste des élèves concernés}},
  coltitle=primary,attach boxed title to top left={yshift=-1.8mm,xshift=3mm},
  boxed title style={colback=primarylight,boxrule=0.4pt,colframe=cardborder,arc=1mm,top=0.5pt,bottom=0.5pt}
]
\vspace{1pt}
{\scriptsize
\begin{tabularx}{\linewidth}{@{}l X l X@{}}
1. & \makebox[\linewidth]{\dotfill} & 6. & \makebox[\linewidth]{\dotfill} \\
2. & \makebox[\linewidth]{\dotfill} & 7. & \makebox[\linewidth]{\dotfill} \\
3. & \makebox[\linewidth]{\dotfill} & 8. & \makebox[\linewidth]{\dotfill} \\
4. & \makebox[\linewidth]{\dotfill} & 9. & \makebox[\linewidth]{\dotfill} \\
5. & \makebox[\linewidth]{\dotfill} & 10. & \makebox[\linewidth]{\dotfill} \\
\end{tabularx}
}
\end{tcolorbox}

\vspace{1pt}
\textbf{\faComments\ I) Mise en situation :}\\
{\small
$\bullet$ Proposer aux élèves de citer quelques noms de métiers.\\
$\bullet$ Que porte un médecin ? \textit{Une blouse blanche (tenue médicale) et du matériel médical.}
}

\vspace{2pt}
\textbf{\faEye\ II) Phase d'identification :}\\
{\small\textit{Proposer aux élèves les trois textes suivants sans titres, puis leur demander de les associer aux titres et images :}}

\vspace{2pt}
\begin{tcolorbox}[
  enhanced,
  colback=blue!3!white,
  colframe=primary!60,
  arc=1.5mm,leftrule=3.5pt,rightrule=0.5pt,toprule=0.5pt,bottomrule=0.5pt,
  top=3pt,bottom=3pt,left=7pt,right=7pt,
  title={\small\textbf{Texte 1}},
  coltitle=primary!90!black,attach boxed title to top left={yshift=-1.8mm,xshift=3mm},
  boxed title style={colback=white,boxrule=0.4pt,colframe=primary!60,arc=1mm,top=0.5pt,bottom=0.5pt}
]
\vspace{1pt}
{\small
Le pompier intervient lors des incendies pour éteindre le feu et sauver les personnes en danger. Il utilise des tuyaux, des lances à eau et porte une tenue spéciale pour se protéger. Le pompier aide aussi lors des accidents et des inondations. Le métier de pompier est très noble car il permet de sauver des vies.
}
\end{tcolorbox}

\vspace{-2pt}
\begin{tcolorbox}[
  enhanced,
  colback=emerald!5!white,
  colframe=secondary,
  arc=1.5mm,leftrule=3.5pt,rightrule=0.5pt,toprule=0.5pt,bottomrule=0.5pt,
  top=3pt,bottom=3pt,left=7pt,right=7pt,
  title={\small\textbf{Texte 2}},
  coltitle=secondary!90!black,attach boxed title to top left={yshift=-1.8mm,xshift=3mm},
  boxed title style={colback=white,boxrule=0.4pt,colframe=secondary,arc=1mm,top=0.5pt,bottom=0.5pt}
]
\vspace{1pt}
{\small
Le lapin est un animal domestique. Il est doux et apprécié des enfants. Il se nourrit de carottes, de légumes et d'herbe. Il vit dans une cage ou dans un clapier et peut être élevé à la ferme ou dans la maison.
}
\end{tcolorbox}

\vspace{-2pt}
\begin{tcolorbox}[
  enhanced,
  colback=amber!5!white,
  colframe=accent,
  arc=1.5mm,leftrule=3.5pt,rightrule=0.5pt,toprule=0.5pt,bottomrule=0.5pt,
  top=3pt,bottom=3pt,left=7pt,right=7pt,
  title={\small\textbf{Texte 3}},
  coltitle=accent!90!black,attach boxed title to top left={yshift=-1.8mm,xshift=3mm},
  boxed title style={colback=white,boxrule=0.4pt,colframe=accent,arc=1mm,top=0.5pt,bottom=0.5pt}
]
\vspace{1pt}
{\small
L'environnement est le lieu qui nous entoure. Pour vivre sainement, nous devons protéger ce milieu de tout ce qui peut le polluer. Pour cela, la participation de tous les citoyens est nécessaire.
}
\end{tcolorbox}

\vspace{1pt}
{\small
\textbf{\faLink\ Consigne : Relie chaque texte à son titre et au mot-clé correspondant :}\\
\vspace{1pt}
\begin{tabularx}{\linewidth}{XXX}
\badge[primary]{* Le pompier} & \badge[secondary]{* Le lapin} & \badge[accent]{* L'environnement} \\
\badge[primary]{* Sauver} & \badge[secondary]{* Domestique} & \badge[accent]{* Polluer} \\
\end{tabularx}
}

\vspace{3pt}
\begin{minipage}{0.32\textwidth}
\centering
\begin{tcolorbox}[enhanced,colback=white,colframe=cardborder,arc=2mm,top=2pt,bottom=2pt,left=2pt,right=2pt]
\centering
\includegraphics[height=3.4cm,keepaspectratio]{images/pompier.png}\\[2pt]
{\scriptsize\textbf{\faFireExtinguisher\ Le Pompier}}
\end{tcolorbox}
\end{minipage}
\hfill
\begin{minipage}{0.32\textwidth}
\centering
\begin{tcolorbox}[enhanced,colback=white,colframe=cardborder,arc=2mm,top=2pt,bottom=2pt,left=2pt,right=2pt]
\centering
\includegraphics[height=3.4cm,keepaspectratio]{images/lapin.png}\\[2pt]
{\scriptsize\textbf{\faPaw\ Le Lapin}}
\end{tcolorbox}
\end{minipage}
\hfill
\begin{minipage}{0.32\textwidth}
\centering
\begin{tcolorbox}[enhanced,colback=white,colframe=cardborder,arc=2mm,top=2pt,bottom=2pt,left=2pt,right=2pt]
\centering
\includegraphics[height=3.4cm,keepaspectratio]{images/environnement.png}\\[2pt]
{\scriptsize\textbf{\faLeaf\ L'Environnement}}
\end{tcolorbox}
\end{minipage}

\newpage
% =========================================================================
% PAGE 5 : Compétence 03 - Suite (Application & Réinvestissement)
% =========================================================================

\vspace*{2pt}
\badge[secondary]{\faTasks\ III) Phase d'application}

\vspace{3pt}
{\small\textbf{Activité 01 : Choisir le titre approprié}\\
\textit{Lis le texte suivant puis réponds aux questions :}}

\begin{tcolorbox}[
  enhanced,
  colback=amber!4!white,
  colframe=accent!80,
  arc=2mm,leftrule=3.5pt,rightrule=0.5pt,toprule=0.5pt,bottomrule=0.5pt,
  top=5pt,bottom=5pt,left=8pt,right=8pt,
  title={\small\textbf{\faBookOpen\ Texte :}},
  coltitle=accent!90!black,attach boxed title to top left={yshift=-2mm,xshift=4mm},
  boxed title style={colback=white,boxrule=0.4pt,colframe=accent!80,arc=1mm,top=1pt,bottom=1pt}
]
\vspace{1pt}
{\small
La girafe est un animal sauvage. Elle est le plus grand animal terrestre. Elle se nourrit de feuilles et de fruits qu'elle cueille en haut des arbres grâce à son long cou. Elle vit dans la savane et se déplace en troupeaux.
}
\end{tcolorbox}

\vspace{3pt}
{\small
\textbf{1/ Quel titre convient le mieux à ce texte ?}\par\vspace{2pt}
\noindent\begin{tabular}{@{\hspace{2mm}}p{4.5cm} p{4.5cm} p{6cm}@{}}
\qcmCircle \textbf{a)} Le lapin & \qcmCircle \textbf{b)} La girafe & \qcmCircle \textbf{c)} Le tremblement de terre \\
\end{tabular}\par\vspace{3pt}

\textbf{2/ Relève du texte le champ lexical relatif au mot \guillemotleft\ Animal\ \guillemotright\ :}\par\vspace{2pt}
\noindent\makebox[\linewidth]{\hspace*{2mm}\textbf{Réponse :} \dotfill}
}

\vspace{4pt}
\badge[primary]{\faPencilRuler\ IV) Phase de réinvestissement}

\vspace{2pt}
{\small
\textbf{Activité 02 : Complète le texte à l'aide de la boîte à mots :}\\
\centerline{\badge[accent]{\textbf{animaux}} \quad \badge[accent]{\textbf{grimaces}} \quad \badge[accent]{\textbf{girafes}} \quad \badge[accent]{\textbf{pain}}}

\vspace{3pt}
\begin{tcolorbox}[enhanced,colback=secondarylight,colframe=secondary,arc=2mm,top=4pt,bottom=4pt,left=8pt,right=8pt]
{\small
J'ai visité le Zoo avec ma famille. J'ai vu beaucoup d'\makebox[2.4cm]{\dotfill}\ . Les singes m'ont beaucoup amusé : ils faisaient des \makebox[2.4cm]{\dotfill}\ . J'ai donné à manger aux \makebox[2.4cm]{\dotfill}\ . \guillemotleft\ Il est interdit de donner du \makebox[2cm]{\dotfill}\ aux animaux !\ \guillemotright\ a rappelé le gardien.
}
\end{tcolorbox}
}

\vspace{4pt}
\begin{minipage}{0.48\textwidth}
\centering
\begin{tcolorbox}[enhanced,colback=white,colframe=cardborder,arc=2.5mm,top=3pt,bottom=3pt,left=3pt,right=3pt]
\centering
\includegraphics[width=\linewidth,height=8.8cm,keepaspectratio]{images/girafe_arbre.png}\\[3pt]
{\scriptsize\textbf{\faTree\ La girafe cueillant les feuilles en haut des arbres}}
\end{tcolorbox}
\end{minipage}
\hfill
\begin{minipage}{0.48\textwidth}
\centering
\begin{tcolorbox}[enhanced,colback=white,colframe=cardborder,arc=2.5mm,top=3pt,bottom=3pt,left=3pt,right=3pt]
\centering
\includegraphics[width=\linewidth,height=8.8cm,keepaspectratio]{images/girafe_savane.png}\\[3pt]
{\scriptsize\textbf{\faSun\ La girafe dans son habitat naturel (la savane)}}
\end{tcolorbox}
\end{minipage}

\newpage
% =========================================================================
% PAGE 6 : Compétence 04 - Production Écrite (Fiche Enseignant)
% =========================================================================
\ficheHeader{1 Heure}{\makebox[2.5cm]{\dotfill}}{1\textsuperscript{ère} AM}{2026 / 2027}

\vspace{-2pt}
\competenceBanner{\faPen\ Compétence 04 : Production Écrite (Le lion)}

\begin{minipage}[t]{0.48\textwidth}
\begin{tcolorbox}[
  enhanced,
  colback=secondarylight,
  colframe=secondary,
  arc=1.5mm,leftrule=3pt,rightrule=0.5pt,toprule=0.5pt,bottomrule=0.5pt,
  top=3pt,bottom=3pt,left=6pt,right=6pt,
  title={\small\textbf{\faBullseye\ Critères d'évaluation}},
  coltitle=secondary!90!black,attach boxed title to top left={yshift=-2mm,xshift=4mm},
  boxed title style={colback=white,boxrule=0.5pt,colframe=secondary,arc=1mm,top=1pt,bottom=1pt}
]
\vspace{1pt}
{\scriptsize
\begin{itemize}[leftmargin=*,noitemsep,topsep=1pt]
  \item[\faCheckCircle] Pertinence du texte produit (respect de la consigne).
  \item[\faCheckCircle] Cohérence textuelle et progression logique.
  \item[\faCheckCircle] Correction de la langue (grammaire, orthographe).
  \item[\faCheckCircle] Lisibilité de l'écrit et qualité de la graphie.
\end{itemize}
}
\end{tcolorbox}
\end{minipage}
\hfill
\begin{minipage}[t]{0.48\textwidth}
\begin{tcolorbox}[
  enhanced,
  colback=accentlight,
  colframe=accent,
  arc=1.5mm,leftrule=3pt,rightrule=0.5pt,toprule=0.5pt,bottomrule=0.5pt,
  top=3pt,bottom=3pt,left=6pt,right=6pt,
  title={\small\textbf{\faListOl\ Consignes de rédaction}},
  coltitle=accent!90!black,attach boxed title to top left={yshift=-2mm,xshift=4mm},
  boxed title style={colback=white,boxrule=0.5pt,colframe=accent,arc=1mm,top=1pt,bottom=1pt}
]
\vspace{1pt}
{\scriptsize
\begin{itemize}[leftmargin=*,noitemsep,topsep=1pt]
  \item[\faPen] Donne un titre accrocheur à ton texte.
  \item[\faFile*] Rédige un court texte de 4 à 5 phrases courtes.
  \item[\faCheck] Utilise le présent de l'indicatif.
  \item[\faExclamation] Ponctue soigneusement ton paragraphe.
\end{itemize}
}
\end{tcolorbox}
\end{minipage}

\vspace{1pt}
\begin{tcolorbox}[
  enhanced,
  colback=white,
  colframe=cardborder,
  arc=1.5mm,
  top=2pt,bottom=2pt,left=6pt,right=6pt,
  title={\scriptsize\textbf{\faUsers\ Élèves concernés par la remédiation en écriture}},
  coltitle=primary,attach boxed title to top left={yshift=-1.8mm,xshift=3mm},
  boxed title style={colback=primarylight,boxrule=0.4pt,colframe=cardborder,arc=1mm,top=0.5pt,bottom=0.5pt}
]
\vspace{1pt}
{\scriptsize
\begin{tabularx}{\linewidth}{@{}l X l X@{}}
1. & \makebox[\linewidth]{\dotfill} & 6. & \makebox[\linewidth]{\dotfill} \\
2. & \makebox[\linewidth]{\dotfill} & 7. & \makebox[\linewidth]{\dotfill} \\
3. & \makebox[\linewidth]{\dotfill} & 8. & \makebox[\linewidth]{\dotfill} \\
4. & \makebox[\linewidth]{\dotfill} & 9. & \makebox[\linewidth]{\dotfill} \\
5. & \makebox[\linewidth]{\dotfill} & 10. & \makebox[\linewidth]{\dotfill} \\
\end{tabularx}
}
\end{tcolorbox}

\vspace{1pt}
\begin{tcolorbox}[
  enhanced,
  colback=amber!4!white,
  colframe=accent!80,
  arc=2mm,
  top=4pt,bottom=4pt,left=8pt,right=8pt,
  title={\small\textbf{\faIdCard\ Fiche d'identité documentaire : Le Lion}},
  coltitle=accent!90!black,attach boxed title to top left={yshift=-2mm,xshift=4mm},
  boxed title style={colback=white,boxrule=0.4pt,colframe=accent!80,arc=1mm,top=1pt,bottom=1pt}
]
\vspace{1pt}
\begin{minipage}[c]{0.58\textwidth}
{\small
\begin{tabular}{@{}ll@{}}
\textbf{\faTag\ Nom :} & Lion \\
\textbf{\faVenus\ La femelle :} & Lionne \\
\textbf{\faPaw\ Le petit :} & Lionceau \\
\textbf{\faMapMarker*\ Habitat :} & Dans la savane d'Afrique \\
\textbf{\faUtensils\ Alimentation :} & De la viande (\textit{carnivore}) \\
\textbf{\faWeight\ Poids :} & Jusqu'à 200 kg \\
\textbf{\faStar\ Aspect :} & Grand, fort et majestueux \\
\end{tabular}
}
\end{minipage}
\hfill
\begin{minipage}[c]{0.38\textwidth}
\centering
\includegraphics[width=\linewidth,height=3.8cm,keepaspectratio]{images/lion_corps.png}
\end{minipage}
\end{tcolorbox}

\vspace{3pt}
\begin{tcolorbox}[
  enhanced,
  colback=white,
  colframe=cardborder,
  arc=2mm,
  top=4pt,bottom=4pt,left=8pt,right=8pt
]
\begin{minipage}[t]{0.63\textwidth}
\vspace{0pt}
\textbf{\faPenSquare\ Espace de Rédaction :}\\[6pt]
\centerline{\textbf{Titre :} \makebox[4.5cm]{\dotfill}}\\[14pt]
\noindent\makebox[\linewidth]{\dotfill}\\[14pt]
\noindent\makebox[\linewidth]{\dotfill}\\[14pt]
\noindent\makebox[\linewidth]{\dotfill}\\[14pt]
\noindent\makebox[\linewidth]{\dotfill}\\[14pt]
\noindent\makebox[\linewidth]{\dotfill}\\[14pt]
\noindent\makebox[\linewidth]{\dotfill}
\end{minipage}
\hfill
\begin{minipage}[t]{0.34\textwidth}
\vspace{0pt}
\centering
\includegraphics[width=\linewidth,height=5.8cm,keepaspectratio]{images/lion_famille.png}
\end{minipage}
\end{tcolorbox}

\newpage
% =========================================================================
% PAGE 7 : Planche de 6 Billets de Production Écrite (Lion) à découper
% =========================================================================

\newcommand{\modernLionTicket}{%
\begin{tcolorbox}[
  enhanced,
  colback=white,
  colframe=primary!70,
  arc=2mm,
  boxrule=0.7pt,
  top=3pt,bottom=3pt,left=5pt,right=5pt,
  height=8.5cm
]
\begin{minipage}[t]{0.58\linewidth}
\textbf{\scriptsize\faPen\ Production écrite : Le lion}\\
{\tiny Nom : \makebox[2.2cm]{\dotfill} Cl. : \makebox[1cm]{\dotfill}}\\[2pt]
{\tiny\textbf{Consigne :}\\
À partir de la fiche ci-dessous, écris un court texte (4 à 5 phrases) pour présenter le lion.
\begin{itemize}[leftmargin=*,noitemsep,topsep=1pt]
  \item Donne un titre à ton texte.
  \item Fais des phrases courtes.
  \item Ponctue soigneusement.
\end{itemize}
}
\vspace{1pt}
\begin{tcolorbox}[enhanced,colback=amber!5,colframe=accent!60,arc=1mm,top=1.5pt,bottom=1.5pt,left=3pt,right=3pt,boxrule=0.4pt]
{\tiny
\textbf{\faIdCard\ Fiche du lion :}\\
$\bullet$ \textbf{Nom :} Lion\\
$\bullet$ \textbf{Femelle :} Lionne\\
$\bullet$ \textbf{Petit :} Lionceau\\
$\bullet$ \textbf{Habitat :} Savane d'Afrique\\
$\bullet$ \textbf{Régime :} Viande (\textit{carnivore})\\
$\bullet$ \textbf{Poids :} Jusqu'à 200 kg\\
$\bullet$ \textbf{Corps :} Grand et fort
}
\end{tcolorbox}
\end{minipage}
\hfill
\begin{minipage}[t]{0.39\linewidth}
\centering
\vspace{10pt}
\includegraphics[width=\linewidth,height=7cm,keepaspectratio]{images/lion_marche.png}
\end{minipage}
\end{tcolorbox}
}

\newcommand{\lionTicketsPage}{%
\noindent
\begin{minipage}[t]{0.48\textwidth}
\modernLionTicket
\vspace{-2pt}
\cutLineH
\vspace{-2pt}
\modernLionTicket
\vspace{-2pt}
\cutLineH
\vspace{-2pt}
\modernLionTicket
\end{minipage}
\hfill
\begin{minipage}[t]{0.48\textwidth}
\modernLionTicket
\vspace{-2pt}
\cutLineH
\vspace{-2pt}
\modernLionTicket
\vspace{-2pt}
\cutLineH
\vspace{-2pt}
\modernLionTicket
\end{minipage}%
}

% Page 7 : Billets Production Écrite (Exemplaire 1/3)
\lionTicketsPage

\newpage
% Page 8 : Billets Production Écrite (Exemplaire 2/3)
\lionTicketsPage

\newpage
% Page 9 : Billets Production Écrite (Exemplaire 3/3)
\lionTicketsPage

\newpage
% =========================================================================
% PAGE 10 : Planche de 15 Vignettes Découpables - Schéma Gazelle
% =========================================================================

\newcommand{\modernGazelleSticker}{%
\begin{tcolorbox}[
  enhanced,
  colback=white,
  colframe=gray!50,
  arc=1.5mm,
  boxrule=0.5pt,
  top=1pt,bottom=1pt,left=1pt,right=1pt,
  height=5.1cm
]
\centering
\includegraphics[width=0.98\linewidth,height=4.8cm,keepaspectratio]{images/gazelle_schema.png}
\end{tcolorbox}
}

\newcommand{\gazelleStickersPage}{%
\noindent
\begin{minipage}[t]{0.315\textwidth}
\modernGazelleSticker
\vspace{-4pt}
\modernGazelleSticker
\vspace{-4pt}
\modernGazelleSticker
\vspace{-4pt}
\modernGazelleSticker
\vspace{-4pt}
\modernGazelleSticker
\end{minipage}
\hfill
\begin{minipage}[t]{0.315\textwidth}
\modernGazelleSticker
\vspace{-4pt}
\modernGazelleSticker
\vspace{-4pt}
\modernGazelleSticker
\vspace{-4pt}
\modernGazelleSticker
\vspace{-4pt}
\modernGazelleSticker
\end{minipage}
\hfill
\begin{minipage}[t]{0.315\textwidth}
\modernGazelleSticker
\vspace{-4pt}
\modernGazelleSticker
\vspace{-4pt}
\modernGazelleSticker
\vspace{-4pt}
\modernGazelleSticker
\vspace{-4pt}
\modernGazelleSticker
\end{minipage}%
}

% Page 10 : Planche Vignettes Gazelle (Exemplaire 1/5)
\gazelleStickersPage

\newpage
% Page 11 : Planche Vignettes Gazelle (Exemplaire 2/5)
\gazelleStickersPage

\newpage
% Page 12 : Planche Vignettes Gazelle (Exemplaire 3/5)
\gazelleStickersPage

\newpage
% Page 13 : Planche Vignettes Gazelle (Exemplaire 4/5)
\gazelleStickersPage

\newpage
% Page 14 : Planche Vignettes Gazelle (Exemplaire 5/5)
\gazelleStickersPage

\end{document}
